\lecture{10}{26 March.}{}
\begin{corollary}
    For any open box 
    \[
        B = \prod _{i=1}^n (a_i, b_i) = \left\{ (x_1, \dots , x_n) \in \mathbb{R} ^n \mid a_i < x_i < b_i \text{ for } 1 \le i \le n  \right\},
    \]
    we have \(m^*(B) = \prod _{i=1}^n (b_i - a_i)\). 
\end{corollary}

\begin{proof}
    First since \(B\) cover itself, we have \(m^*(B) \le \mathrm{vol}(B) = \prod _{i=1}^n (b_i - a_i) \). Choose \(\varepsilon_i < \frac{b_i - a_i}{2}\). Then, 
    \[
        [a_i + \varepsilon_i, b_i - \varepsilon_i] \subseteq (a_i, b_i).
    \]   
    Hence,
    \[
        \prod _{i=1}^n [a_i + \varepsilon_i, b_i - \varepsilon_i] \subseteq \prod _{i=1}^n (a_i, b_i),
    \]
    which shows 
    \[
        m^* \left( \prod _{i=1}^n [a_i + \varepsilon_i, b_i - \varepsilon_i] \right) \le m^* (B), 
    \]
    and by previous proposition, we have 
    \[
        \prod _{i=1}^n (b_i - a_i - 2 \varepsilon_i) = m^* \left( \prod _{i=1}^n [a_i + \varepsilon_i, b_i - \varepsilon_i] \right) \le m^* (B).
    \]
    Now pick all \(\varepsilon _i \to 0\), then we have \(m^*(B) \ge \mathrm{vol} (B) \).  
\end{proof}

\begin{eg}
    Let \(r > 0\), then \((-r, r) \subseteq \mathbb{R} \), and thus 
    \[
        2r = m^* ((-r, r)) \le m^* (\mathbb{R} ) 
    \]  
    for all \(r\), so \(m^*(\mathbb{R} ) = \infty \). Similarly, \(m^*(\mathbb{R} ^n) = \infty \).   
\end{eg}

\begin{eg}
    \[
        \mathbb{Q} = \text{ the set of rational numbers } = \bigcup_{q \in \mathbb{Q} } \left\{ q \right\},   
    \]
    and \(\mathbb{Q} \) is countable, so 
    \[
        m^*(\mathbb{Q} ) \le \sum_{q \in \mathbb{Q} } m^*(\left\{ q \right\} ) = 0. 
    \]
    In particular, \(\exists \) an countable family of open boxes \((B_j)\) s.t. \(\mathbb{Q} \subseteq \bigcup_{j \in J} B_j \) and 
    \[
        \sum_{j \in J} \mathrm{vol}(B_j) \le m^*(\mathbb{Q} ) + \varepsilon = \varepsilon.  
    \]   
\end{eg}

\begin{eg}
    Since \(\mathbb{R} = (\mathbb{R} \setminus \mathbb{Q} ) \cup \mathbb{Q} \), we have 
    \[
        m^*(\mathbb{R} ) \le m^*(\mathbb{R} \setminus \mathbb{Q} ) + m^*(\mathbb{Q} ),
    \] 
    and since \(m^*(\mathbb{R} ) = \infty\) and \(m^*(\mathbb{Q} ) = 0\), we know \(m^*(\mathbb{R} \setminus \mathbb{Q} ) = \infty \).   
\end{eg}

\begin{claim}
    \(m^* ([0, 1] \setminus \mathbb{Q} ) = 1\). 
\end{claim}
\begin{explanation}
    Since 
    \[
        [0, 1] = ([0, 1] \setminus \mathbb{Q} ) \cup ([0, 1] \cap \mathbb{Q}),
    \]
    we have 
    \[
        1 = m^*([0, 1]) \le m^*([0, 1] \setminus \mathbb{Q} ) + m^* ([0, 1] \cap \mathbb{Q} ) = m^*([0, 1] \setminus \mathbb{Q} )
    \]
    since 
    \[
        [0, 1] \cap \mathbb{Q} \subseteq \mathbb{Q} \text{ and } 0 \le m^*([0, 1] \cap \mathbb{Q} ) \le m^*(\mathbb{Q} ) = 0. 
    \]
    Now since 
    \[
        [0, 1] \setminus \mathbb{Q} \subseteq [0, 1],
    \]
    so 
    \[
        m^*([0, 1] \setminus \mathbb{Q} ) \le m^* ([0, 1]) = 1.
    \]
    Thus,
    \[
        m^*([0, 1] \setminus \mathbb{Q} ) = 1.
    \]
\end{explanation}

\begin{claim}
    Let 
    \[
        L = \left\{ (x, 0) \mid x \in \mathbb{R} \right\} \subseteq \mathbb{R}^2,
    \]
    then \(m^*(L) = 0\). 
\end{claim}

\begin{explanation}
    Let \(L_n = \left\{ (x, 0) \mid -n \le x \le n \right\} \subseteq \mathbb{R} ^2 \), then 
    \[
        L_n \subseteq [-n, n] \times [-\varepsilon, \varepsilon],
    \]
    so 
    \[
        0 \le m^*(L_n) \le (2n) \cdot (2 \varepsilon).
    \]
    Let \(\varepsilon \to 0\), then we have 
    \[
        m^*(L_n) = 0.
    \] 
    Now since 
    \[
        L = \bigcup_{n=1}^{\infty} L_n, 
    \]
    so 
    \[
        m^*(L) \le \sum_{n=1}^{\infty} m^*(L_n) = 0, 
    \]
    which shows \(m^*(L) = 0\). 
\end{explanation}

\section{Outer measure is not additive}
We only know that
\[
    m^* \left( \bigcup_{j \in J} A_j \right) \le \sum_{j \in J} m^*(A_j)  
\]
if \(J\) is countable. To construct a measure, we require that if \(A_i \cap A_j = \varnothing\) for all \(i, j \in J\) and \(i \neq j\),
\[
    m^* \left( \bigcup_{j \in J} A_j \right) = \sum_{j \in J} m^*(A_j). 
\]    

\begin{proposition}[Failure of countable additivity]
    There exists a countable collection \((A_j)_{j \in J}\) of disjoint subsets of \(\mathbb{R} \) s.t. 
    \[
        m^* \left( \bigcup_{j \in J} A_j  \right) \neq \sum_{j \in J} m^*(A_j). 
    \]  
\end{proposition}

\begin{proof}
    We shall construct such a family from the coset of \(\mathbb{Q} \) in \(\mathbb{R} \). 
    \begin{itemize}
        \item Step 1. Coset of \(\mathbb{Q} \). For \(x \in \mathbb{R} \), let 
        \[
            x + \mathbb{Q} = \left\{ x + q : q \in \mathbb{Q} \right\}. 
        \]
        This is called a coset of \(\mathbb{Q} \). Two properties of cosets:
        \begin{itemize}
            \item [(1)] \((x + \mathbb{Q} ) \cap (y + \mathbb{Q} ) \neq \varnothing\) implies \(x + \mathbb{Q} = y + \mathbb{Q} \).  
            \item [(2)] Every coset meets the interval \([0, 1]\). (We can pick \(-x \le q \le 1 - x\) so that \(0 \le q + x \le 1\)) 
        \end{itemize} 
        \item Step 2. Choose one representative from each coset. Let \(\mathbb{R} / \mathbb{Q} \) denote the collection of all cosets. For each coset \(A \in \mathbb{R} / \mathbb{Q} \). By (2) in Step 1, there exists \(x_A \in [0, 1]\) s.t. \(x_A \in A \cap [0, 1]\)(By axiom of choice), we can choose one and define 
        \[
            E = \left\{ x_A : A \in \mathbb{R} / \mathbb{Q} , x_A \in [0, 1] \right\} \subseteq [0, 1]. 
        \]
        \item Step 3. Rational translates of \(E\). Consider the set 
        \[
            X \coloneqq \bigcup_{q \in \mathbb{Q} \cap [-1, 1]} (q + E). 
        \]
        Since \(E \subseteq [0, 1]\) and \(-1 \le q \le 1\), we have \(X \subseteq [-1, 2]\). Now we claim \([0, 1] \subseteq X\). Let \(y\) be any number in \([0, 1]\). Let \(y \in [0, 1]\), then \(y\) is in some coset of \(\mathbb{R} / \mathbb{Q} \). Thus, \(\exists \) cosets \(A\) and \(x_A \in E \subseteq [0, 1]\) s.t. \(y \in A = x_A + \mathbb{Q} = y + \mathbb{Q} \). Let \(y = x_A + q\) for some \(q \in \mathbb{Q}\). Thus, \(q = y - x_A \in [-1, 1]\). Thus, \(q \in \mathbb{Q} \cap [-1, 1]\). Hence,
        \[
            y = q + x_A \in q + E \subseteq X.
        \]               
        Thus, 
        \[
            [0, 1] \subseteq X \subseteq [-1, 2].
        \]
        \item Step 4. the sets \(q + E\) are pairwise disjoint. We now show that the family 
        \[
            \left\{ q + E  \right\}_{q \in \mathbb{Q} \cap [-1, 1]} 
        \]
        is pairwise disjoint. Suppose, for contradiction, there are distinct \(q, q^{\prime} \in \mathbb{Q} \cap [-1, 1]\) we had
        \[
            (q + E) \cap (q^{\prime} + E) \neq \varnothing.
        \] 
        Then, there exists cosets \(A, A^{\prime} \in \mathbb{R} / \mathbb{Q} \) s.t. 
        \[
            q + x_A = q^{\prime} + x_{A^{\prime} }.
        \] 
        Hence,
        \[
             x_{A^{\prime} } - x_A = q - q^{\prime} \in \mathbb{Q}.
        \]
        Now since \(x_A \in A = x_A + \mathbb{Q} \) and \(x_{A^{\prime} } \in A^{\prime} = x_{A^{\prime} } + \mathbb{Q} \), so \(A = A^{\prime}\). But \(E\) contains exactly one representative from each coset, which means \(x_A = x_{A^{\prime}}\). Thus, \(q = q^{\prime}\), which contradicts to the assumption, so the sets \(q + E\) are pairwise disjoint.
        \item Step 5. Comparing the two sides. Now since \([0, 1] \subseteq X \subseteq [-1, 2]\), 
        \[
            1 = m^*([0, 1]) \le m^*(X) \le m^*([-1, 2]) = 3.
        \]
        On the other hand, 
        \[
            m^*(q + E) = m^*(E) \text{ for every } q \in \mathbb{Q} \cap [-1, 1]. 
        \]
        Thus,
        \[
            \sum_{q \in \mathbb{Q} \cap [-1, 1]} m^*(q + E) = \sum_{q \in \mathbb{Q} \cap [-1, 1]} m^*(E).  
        \]
        Note that 
        \[
            \sum_{q \in \mathbb{Q} \cap [-1, 1]} m^*(E) = \begin{dcases}
                0, &\text{ if } m^*(E) = 0;\\
                \infty , &\text{ if } m^*(E) > 0.
            \end{dcases} 
        \]
        In particular, it cannot lie in the interval \([1, 3]\). But \(m^*(X) \in [1, 3]\). Hence,
        \[
            m^*(X) \neq \sum_{q \in \mathbb{Q} \cap [-1, 1]} m^*(q + E).
        \]  
        Since \(\mathbb{Q} \cap [-1, 1]\) is countably infinite, so we're done. 
    \end{itemize}  
\end{proof}

\begin{proposition}
    There exists a finite collection \((A_j)_{j \in J}\) of disjoint subsets of \(\mathbb{R} \) s.t. 
    \[
        m^* \left( \bigcup_{j \in J} A_j  \right) \neq \sum_{j \in J} m^*(A_j). 
    \]  
\end{proposition}

\begin{proof}
    Suppose \(m^*\) was finitely additive, and let \(E\) and \(X\) to be same as previous proposition, then 
    \[
        m^*(X) \le \sum_{q \in \mathbb{Q} \cap [-1, 1]} m^*(q + E) = \sum_{q \in \mathbb{Q} \cap [-1, 1]} m^*(E).   
    \]  
    Since we know that \(1 \le m^*(X) \le 3\), we thus have \(m^*(E) \neq 0\), or otherwise we will have \(m^*(X) \le 0\), but we have \(m^*(X) \in [1, 3]\). 

    Since \(m^*(E) > 0\), so \(m^*(E) > \frac{1}{n}\) for some \(n \in \mathbb{N}\). Now let \(J\) be a finite subset of \(\mathbb{Q} \cap [-1, 1]\) of cardinality \(3n\). If \(m^*\) were finitely additive, then we would have 
    \[
        m^* \left( \bigcup_{j \in J} (q + E)  \right) = \sum_{q \in J} m^*(q + E) = \sum_{q \in J} m^*(E) > 3n \cdot \frac{1}{n} = 3. 
    \] 
    But we know \(\bigcup_{j \in J} (q + E) \) is a subset of \(X\), so it has outer measure at most \(3\), so we have a contradiction.             
\end{proof}